% Options for packages loaded elsewhere
\PassOptionsToPackage{unicode}{hyperref}
\PassOptionsToPackage{hyphens}{url}
\documentclass[
]{article}
\usepackage{xcolor}
\usepackage{amsmath,amssymb}
\setcounter{secnumdepth}{-\maxdimen} % remove section numbering
\usepackage{iftex}
\ifPDFTeX
  \usepackage[T1]{fontenc}
  \usepackage[utf8]{inputenc}
  \usepackage{textcomp} % provide euro and other symbols
\else % if luatex or xetex
  \usepackage{unicode-math} % this also loads fontspec
  \defaultfontfeatures{Scale=MatchLowercase}
  \defaultfontfeatures[\rmfamily]{Ligatures=TeX,Scale=1}
\fi
\usepackage{lmodern}
\ifPDFTeX\else
  % xetex/luatex font selection
\fi
% Use upquote if available, for straight quotes in verbatim environments
\IfFileExists{upquote.sty}{\usepackage{upquote}}{}
\IfFileExists{microtype.sty}{% use microtype if available
  \usepackage[]{microtype}
  \UseMicrotypeSet[protrusion]{basicmath} % disable protrusion for tt fonts
}{}
\makeatletter
\@ifundefined{KOMAClassName}{% if non-KOMA class
  \IfFileExists{parskip.sty}{%
    \usepackage{parskip}
  }{% else
    \setlength{\parindent}{0pt}
    \setlength{\parskip}{6pt plus 2pt minus 1pt}}
}{% if KOMA class
  \KOMAoptions{parskip=half}}
\makeatother
\usepackage{longtable,booktabs,array}
\newcounter{none} % for unnumbered tables
\usepackage{calc} % for calculating minipage widths
% Correct order of tables after \paragraph or \subparagraph
\usepackage{etoolbox}
\makeatletter
\patchcmd\longtable{\par}{\if@noskipsec\mbox{}\fi\par}{}{}
\makeatother
% Allow footnotes in longtable head/foot
\IfFileExists{footnotehyper.sty}{\usepackage{footnotehyper}}{\usepackage{footnote}}
\makesavenoteenv{longtable}
\setlength{\emergencystretch}{3em} % prevent overfull lines
\providecommand{\tightlist}{%
  \setlength{\itemsep}{0pt}\setlength{\parskip}{0pt}}
\usepackage{bookmark}
\IfFileExists{xurl.sty}{\usepackage{xurl}}{} % add URL line breaks if available
\urlstyle{same}
\hypersetup{
  hidelinks,
  pdfcreator={LaTeX via pandoc}}

\author{}
\date{}

\begin{document}

\section{On the Ontological Nature of Mass and Energy: A Chronometric
Monism from the Theory of Objective
Time}\label{on-the-ontological-nature-of-mass-and-energy-a-chronometric-monism-from-the-theory-of-objective-time}

\textbf{Luís Pimenta Damásio}\\
Independent Researcher (Philosophy of Time \& Foundations of Physics)
--- Lisbon, Portugal\\
ORCID:
\href{https://orcid.org/0009-0001-7606-7557}{0009-0001-7606-7557}\\
Master Compendium Reference DOI: \texttt{10.5281/zenodo.24163322}

\begin{center}\rule{0.5\linewidth}{0.5pt}\end{center}

\subsection{Abstract}\label{abstract}

Modern physics possesses extraordinarily precise mathematical formalisms
for calculating how mass, energy, frequency, and fields interact, yet it
remains ontologically silent regarding what these entities \emph{are} in
themselves. Textbooks routinely define energy circularly as the capacity
to do work and mass operationally as resistance to acceleration or
coupling to an underlying field, leaving fundamental metaphysics trapped
in a dualism between inert material substance (\(m\)) and dynamic
activity (\(E\)) unfolding inside a passive coordinate container
(\(t\)). Drawing upon the Theory of Objective Time
(\(\text{AEON-}T_{\mu\nu}\)), this treatise develops a systematic
chronometric monism in which \textbf{Objective Time}---formalized as
\textbf{Temporal Intensity (\(\rho_\tau\))}, the active rate at which
un-condensed relational potentiality (\(\Omega\)) resolves into
actualized history (\(v \to v + 1\))---serves as the sole ontological
primitive. By reifying Kant's transcendental aesthetic into the
constitutive structure of reality itself and uniting it with Einstein's
relativistic clock metrology via Two-Stage Condensation (\(\Phi_\tau\)),
the paper strips away anthropogenic unit-conversion scaffolding
(\(h, c^2, k_B, G\)) and reformulates \(E = h\nu = m_0 c^2\) into a
single ontological conservation and allocation principle: the
\textbf{Cadence Partition Law}
(\(\rho_{\tau,\text{total}}^2 = \rho_{\tau,\text{int}}^2 + \rho_{\tau,\text{kin}}^2\)).
Within this framework, energy is total resolution throughput, rest mass
is topologically folded, self-referential internal cadence
(\(\rho_{\tau,\text{int}}\), forming a localized \emph{Center of
Resolution}), inertia and gravity emerge as the internal phase-rigidity
and external pacing refraction (\(n_{\text{eff}}\)) of that same
standing lock, and quantum fields are revealed as un-condensed
topological modalities of the Functional Infinity (\(\Omega\)). While
articulated strictly as a treatise in fundamental ontology, its
categories project seamlessly onto modern physics as a complete,
non-dualistic ontology of matter, radiation, spacetime, and complexity.

\begin{center}\rule{0.5\linewidth}{0.5pt}\end{center}

\subsection{1. The Ontological Silence of Modern
Physics}\label{the-ontological-silence-of-modern-physics}

A persistent divide separates mathematical physics from fundamental
ontology. Physics excels at specifying functional relationships among
measured quantities---how a trajectory bends, how a spectral line
shifts, how a cross-section scales---yet when pressed on the intrinsic
nature of its most elementary referents, its definitions become
conspicuously circular.

Consider the foundational triad of mechanics and relativity:
\emph{mass}, \emph{energy}, and \emph{time}. In standard pedagogical and
foundational treatments, energy is defined as ``the capacity to do
work,'' while work is defined as force integrated over displacement,
force as the rate of change of momentum, and momentum as mass multiplied
by velocity---where velocity itself presupposes a pre-existing,
background coordinate parameter \(t\) against which spatial change is
measured. Mass, in turn, is characterized either through its dynamical
effects (inertial resistance \(F/a\) or gravitational coupling) or as an
intrinsic numerical parameter assigned to fields in a Lagrangian. Even
Einstein's celebrated equivalence,

\[E = h\nu = m_0 c^2,\]

is almost universally interpreted as a quantitative conversion rule
linking three distinct ontological categories: \emph{energy} (dynamic
activity), \emph{frequency} (periodic rate), and \emph{mass} (material
substance).

For the metaphysician, this leaves the central riddle untouched. If mass
is conceived as a passive, persistent substrate (\emph{substance}) and
energy as an active capacity or flux (\emph{process}), how can inert
stuff be ``converted'' into pure activity during nuclear fission or
particle--antiparticle annihilation? Conversely, if everything is
reduced to pure mathematical formalism without an ontological referent,
physics degenerates into instrumentalist bookkeeping---predicting
instrument readings while forfeiting any account of what exists in
reality.

This treatise argues that the impasse arises from a single inverted
premise: \textbf{treating time as a passive, empty container (\(t\))
within which substances persist and energies flow, rather than
recognizing Objective Time as the active ontological ground of existence
itself.}

\begin{center}\rule{0.5\linewidth}{0.5pt}\end{center}

\subsection{\texorpdfstring{2. Objective Time (\(\rho_\tau\)),
Subjective Clocks, and the Kant--Einstein
Synthesis}{2. Objective Time (\textbackslash rho\_\textbackslash tau), Subjective Clocks, and the Kant--Einstein Synthesis}}\label{objective-time-rho_tau-subjective-clocks-and-the-kanteinstein-synthesis}

To resolve the dualism between static substance (\emph{Being}) and
dynamic flux (\emph{Becoming}), ontology must unite two insights that
modern thought has kept artificially separated: Immanuel Kant's
transcendental insight into the constitutive primacy of time, and Albert
Einstein's operational insight into how physical clocks behave \emph{in
situ}.

In the \emph{Critique of Pure Reason}, Kant recognized that time and
space are not empirical objects floating inside the world; they are the
\emph{a priori} constitutive conditions that make sequence, persistence,
and relationality possible---a \textbf{three-way infinity} (endless
past, endless future, and infinitely divisible depth). Yet to avoid
substantivalist antinomies, Kant confined that constitutive architecture
strictly to the transcendental subject
(\emph{objects-as-experienced-by-us}), leaving the
\emph{thing-in-itself} (\emph{object-as-object}) outside time and space.
Einstein, conversely, proved experimentally that what clocks measure
\emph{in situ} is local, relational, and dynamically coupled to motion
and gravitating mass (\(d\tau = \sqrt{N^2}\,dt\)). Yet by flattening
time into a static four-dimensional geometric manifold, physics
spatialized becoming into a frozen block universe.

The Theory of Objective Time (\(\text{AEON-}T_{\mu\nu}\)) heals this
rift by \textbf{reifying Kant's transcendental aesthetic directly into
the constitutive architecture of reality itself}, bridged by two
fundamental modalities:

\begin{enumerate}
\def\labelenumi{\arabic{enumi}.}
\tightlist
\item
  \textbf{The Functional Infinity (\(\Omega\)):} The global reservoir of
  un-condensed, coherent relational potentiality---preserving Kant's
  unbroken three-way infinity as open functional possibility prior to
  discrete actualization, without generating substantival antinomies.
\item
  \textbf{Two-Stage Condensation (\(\Phi_\tau\)):} The ontological act
  by which un-condensed relational potentiality (\(\Omega\)) resolves
  into discrete, finite, irreversible actualizations (\(v \to v + 1\)),
  bounded spatially by the \textbf{Operational Simultaneity Horizon
  (\(\mathcal{L}_{\text{sim}}\))} and governed by the Weyl-invariant
  horizon constant:
\end{enumerate}

\[\mathcal{C}_\tau = \frac{k_B T_{\text{coh}} \cdot \mathcal{L}_{\text{sim}}}{h c} = \frac{1}{2\pi} \quad \Longleftrightarrow \quad \frac{k_B T_{\text{coh}} \cdot \mathcal{L}_{\text{sim}}}{\hbar c} = 1\text{ rad}.\]

From this bridge emerges the sole quantitative primitive of ontology:
\textbf{Temporal Intensity (\(\rho_\tau\))}.

\begin{quote}
\textbf{Axiom 1 (Temporal Intensity as Ontological Primitive):}\\
\emph{Temporal Intensity (\(\rho_\tau\)) is the intrinsic, local rate at
which un-condensed potentiality (\(\Omega\)) resolves into actualized
history (\(v \to v + 1\)) for an object-as-object. Nothing exists prior
to or independent of its cadence of actualization; to be is to resolve
at a determinate intensity \(\rho_\tau\).}
\end{quote}

\subsubsection{\texorpdfstring{Objective Time vs.~Clock Time
(\emph{Object-as-Object}
vs.~\emph{Object-as-Subject})}{Objective Time vs.~Clock Time (Object-as-Object vs.~Object-as-Subject)}}\label{objective-time-vs.-clock-time-object-as-object-vs.-object-as-subject}

Under Axiom 1, \textbf{frequency (\(\nu\)) is not a derivative ratio}
obtained by counting cycles against an external cosmic clock; an
entity's intrinsic oscillation cadence \emph{is} its primary expression
of Objective Time (\(\rho_\tau \equiv \nu\)). What physics and human
experience call ``time'' as measured by a clock is \textbf{Subjective
(Operational) Time}: * A \textbf{clock} is an
\emph{object-as-subject}---a stabilized, countable \textbf{Center of
Resolution} benchmarking external cadences against its own internal
rhythm via transitive signals across a shared \(\Omega\) horizon:
\[\frac{\Delta\tau_A}{\Delta\tau_B} = \frac{\rho_{\tau,A}}{\rho_{\tau,B}}.\]
* Relativistic time dilation does not occur because an empty geometric
container warps, but because one Center of Resolution physically
executes fewer internal actualization cycles per cycle of another.

\subsubsection{The Anthropogenic Scaffolding of Fundamental
Constants}\label{the-anthropogenic-scaffolding-of-fundamental-constants}

Once Temporal Intensity (\(\rho_\tau\)) is recognized as the single
ontological currency of reality, the canonical constants \(h, c, k_B\),
and \(G\) are revealed as \textbf{anthropogenic unit-conversion exchange
rates} translating one single cadence across human laboratory scales: *
\textbf{\(h\)} translates raw resolution cadence (\(\text{Hz}\)) into
mechanical energy (\(\text{J}\)). * \textbf{\(c\)} defines the
un-condensed propagation ratio linking spatial extension (\(\text{m}\))
to substrate resolution cadence (\(\text{s}\)). * \textbf{\(h/c^2\)}
translates self-anchored internal cadence (\(\text{Hz}\)) into inertial
mass (\(\text{kg}\)). * \textbf{\(k_B\)} translates stochastic cadence
dispersion (thermal jitter, \(\text{K}\)) into energy units
(\(\text{J}\)).

\begin{center}\rule{0.5\linewidth}{0.5pt}\end{center}

\subsection{\texorpdfstring{3. Deconstructing \(E = h\nu = m_0 c^2\):
The Cadence Partition
Law}{3. Deconstructing E = h\textbackslash nu = m\_0 c\^{}2: The Cadence Partition Law}}\label{deconstructing-e-hnu-m_0-c2-the-cadence-partition-law}

When we strip away the unit-conversion constants \(h\) and \(c^2\) by
working in native chronometric units (\(h = 1, c = 1\)), Einstein's
triple equivalence collapses into a single ontological conservation and
allocation principle: the \textbf{Cadence Partition Law}.

\[\rho_{\tau,\text{total}}^2 = \underbrace{\rho_{\tau,\text{int}}^2}_{\substack{\text{Immanent Act} \\ (\text{Persistence / Rest Mass } m_0)}} + \underbrace{\rho_{\tau,\text{kin}}^2}_{\substack{\text{Transitive Act} \\ (\text{Relation / Momentum } p)}}\]

Every existent possesses a unified budget of Temporal Intensity
(\(\rho_{\tau,\text{total}} \equiv E/h\)) that partitions orthogonally
between two topological modes: 1. \textbf{Immanent Cadence
(\(\rho_{\tau,\text{int}} \equiv m_0 c^2 / h\)):} Temporal Intensity
folded inward upon itself (via chiral phase-locking,
\(\psi_L \leftrightarrow \psi_R\)) into a closed, self-referential
standing oscillation---a localized \textbf{Center of Resolution} that
sustains self-identical persistence across ticks. 2. \textbf{Transitive
Cadence (\(\rho_{\tau,\text{kin}} \equiv pc / h\)):} Temporal Intensity
directed outward across the relational substrate, shifting an entity's
spatial phase-relation relative to other Centers of Resolution.

Parameterizing this orthogonal partition by the \textbf{Cadence
Partition Angle \(\theta\)} yields:

\[\cos\theta = \frac{\rho_{\tau,\text{int}}}{\rho_{\tau,\text{total}}} = \sqrt{1 - \beta^2} = \frac{1}{\gamma}, \qquad \sin\theta = \frac{\rho_{\tau,\text{kin}}}{\rho_{\tau,\text{total}}} = \beta = \frac{v}{c}.\]

This geometry divides all possible existents into three exhaustive
ontological regimes: * \textbf{Pure Transitive Act
(\(\theta = 90^\circ, \rho_{\tau,\text{int}} = 0\)):} Massless bosons
(photons). Zero internal clock-lock, zero rest frame, zero inertia;
100\% of Temporal Intensity propagates across \(\Omega\) at \(c\). *
\textbf{Pure Immanent Act
(\(\theta = 0^\circ, \rho_{\tau,\text{kin}} = 0\)):} Stationary rest
mass (\(m_0\)). 100\% of Temporal Intensity is locked into internal
self-maintenance. * \textbf{Partitioned Act
(\(0^\circ < \theta < 90^\circ\)):} Moving matter. Tilting \(\theta\)
toward translation reduces the internal self-maintenance projection by
\(\cos\theta = 1/\gamma\), directly explaining kinematic time dilation
and why no entity with \(\rho_{\tau,\text{int}} > 0\) can ever reach
\(\theta = 90^\circ\) (\(c\)) with finite \(\rho_{\tau,\text{total}}\).
During annihilation (\(e^- + e^+ \to 2\gamma\)), no ``inert stuff'' is
destroyed; two closed chiral phase-locks simply unfold from
\(\theta = 0^\circ\) into open transitive propagation at
\(\theta = 90^\circ\).

\begin{center}\rule{0.5\linewidth}{0.5pt}\end{center}

\subsection{4. What Energy, Mass, Inertia, and Gravity Are: The Sunlit
Desk}\label{what-energy-mass-inertia-and-gravity-are-the-sunlit-desk}

Within this chronometric monism, the classical riddles of mechanics
dissolve: * \textbf{Energy} is unrestricted resolution throughput
(\(\rho_{\tau,\text{total}}\))---the total rate of ontological
actualization regardless of topological folding. * \textbf{Rest Mass} is
self-anchored, standing Temporal Intensity
(\(m_0 = (h/c^2)\rho_{\tau,\text{int}}\))---persistent \emph{Being}
revealed as rhythmic, self-referential \emph{Becoming}. *
\textbf{Inertia} is the \textbf{internal phase-rigidity} of a standing
chiral oscillator (\(\rho_{\tau,\text{int}}\)) resisting any external
force attempting to tilt its Cadence Partition angle \(\theta\). *
\textbf{Gravity} is the \textbf{external pacing depletion} induced in
the surrounding substrate by anchoring that same standing cadence
(\(\rho_{\tau,\text{int}}\)), expressed macroscopically as an effective
refractive pacing index \(n_{\text{eff}}(r) = 1/\sqrt{N^2(r)} > 1\) that
refracts neighboring wavepackets inward without spin-2 gravitons.
Because inertial resistance and gravitational pacing refraction both
measure the exact same quantity (\(\rho_{\tau,\text{int}} \ge 0\)), the
Equivalence Principle (\(m_i \equiv m_g\)) and universal gravitational
attraction follow as logical necessities.

\subsubsection{Everyday Illustration: A Tungsten Paperweight on a Sunlit
Desk}\label{everyday-illustration-a-tungsten-paperweight-on-a-sunlit-desk}

Consider a small cube of tungsten sitting on a wooden desk in a sunlit
room. Standard physics struggles to reconcile the unyielding heaviness
of the metal with the fact that atoms are overwhelmingly empty space,
surrounded by invisible air and traversed by light. Under the Theory of
Objective Time, all three elements on the desk are manifestations of the
exact same primitive (\(\rho_\tau\)) organized into distinct topological
regimes: 1. \textbf{The Tungsten Paperweight (Hyper-Dense Immanent
\(\rho_{\tau,\text{int}}\)):} Each tungsten nucleus packs 184 nucleons
cycling as standing chiral oscillators at
\(\rho_{\tau,\text{int}} \sim 2.27 \times 10^{23}\text{ Hz}\). Lifting
the cube requires tilting the partition angle \(\theta\) of
\(\sim 10^{25}\) tightly locked clocks simultaneously
(heaviness/inertia), while Pauli phase-exclusion between standing
fermionic electron shells prevents your fingers from passing through the
metal (solidity). 2. \textbf{The Surrounding Air (Sparse, High-Jitter
\(\rho_\tau\)):} Nitrogen and oxygen molecules carry the same internal
nucleon clocks at \(\sim 15,000\times\) lower spatial density, while
room-temperature thermal jitter
(\(T_{\text{coh}} \approx 295\text{ K}\),
\(\mathcal{L}_{\text{sim}} \approx 1.23\,\mu\text{m}\)) continuously
triggers random collisional condensations (\(\Phi_\tau\)), preventing
lattice lock and letting your hand pass through effortlessly. 3.
\textbf{The Sunlight and Metallic Gleam (Transitive
\(\rho_{\tau,\text{kin}}\) Across \(\Omega\)):} Sunlight enters the room
as pure transitive cadence
(\(\rho_{\tau,\text{kin}} \sim 5 \times 10^{14}\text{ Hz}, \rho_{\tau,\text{int}} = 0\))
propagating across the room's QED relational horizon (\(\Omega\)). It
passes un-condensed through transparent air molecules, couples
coherently to the delocalized conduction-electron phase-sea of the
tungsten surface, and re-radiates toward your eye as a metallic gleam.

\begin{center}\rule{0.5\linewidth}{0.5pt}\end{center}

\subsection{\texorpdfstring{5. Quantum Fields (\(\Omega\)) and the
Cross-Scale Triadic Law
(\(S_{\text{max}} \leftrightarrow \rho_\tau \leftrightarrow K\))}{5. Quantum Fields (\textbackslash Omega) and the Cross-Scale Triadic Law (S\_\{\textbackslash text\{max\}\} \textbackslash leftrightarrow \textbackslash rho\_\textbackslash tau \textbackslash leftrightarrow K)}}\label{quantum-fields-omega-and-the-cross-scale-triadic-law-s_textmax-leftrightarrow-rho_tau-leftrightarrow-k}

Projecting this ontology across modern quantum field theory and
cosmology yields a complete unified picture of physical reality:

\subsubsection{\texorpdfstring{Quantum Fields as Un-Condensed Modalities
of
\(\Omega\)}{Quantum Fields as Un-Condensed Modalities of \textbackslash Omega}}\label{quantum-fields-as-un-condensed-modalities-of-omega}

A quantum field is not an operator-valued fluid filling a pre-existing
spacetime box; it is a \textbf{topological phase-symmetry of the
Functional Infinity (\(\Omega\))} mapping all un-condensed transition
possibilities prior to \(\Phi_\tau\) condensation. Particles are
localized actualizations resolving out of \(\Omega\) (transitive
\(\rho_{\tau,\text{kin}}\) packets for bosons; standing chiral
\(\rho_{\tau,\text{int}}\) oscillators for fermions). Every electron in
the cosmos is identical because each is a standing resonance of the
exact same universal chiral symmetry---not an individual lump of
substance. Moreover, because un-condensed vacuum modes carry zero
localized clock-anchoring
(\(\rho_{\tau,\text{int}} = 0, \nabla n_{\text{eff}} = 0\)), the uniform
vacuum floor induces zero gravitational curvature, dissolving the
\(10^{120}\) cosmological constant catastrophe at its ontological root.

\subsubsection{\texorpdfstring{The Cross-Scale Triadic Law
(\(S_{\text{max}} \leftrightarrow \rho_\tau \leftrightarrow K\))}{The Cross-Scale Triadic Law (S\_\{\textbackslash text\{max\}\} \textbackslash leftrightarrow \textbackslash rho\_\textbackslash tau \textbackslash leftrightarrow K)}}\label{the-cross-scale-triadic-law-s_textmax-leftrightarrow-rho_tau-leftrightarrow-k}

Across all 61 orders of magnitude from the Planck scale
(\(10^{-35}\text{ m}\)) to the cosmic horizon (\(10^{26}\text{ m}\)),
reality obeys a triadic coordination law linking \textbf{Maximal
Potential Branching Entropy (\(S_{\text{max}}\) of \(\Omega\))},
\textbf{Temporal Intensity (\(\rho_\tau\))}, and \textbf{Actualized
Structural Complexity (\(K\))}: * \textbf{The Quantum Chiral Spike
(\(10^{-18}\text{--}10^{-13}\text{ m}\)):} Where chiral symmetry
breaking forges standing Compton clocks, the local density of accessible
quantum degrees of freedom (\(S_{\text{max}}\)) reaches its
peak---requiring a sharp spike in Temporal Intensity
(\(\rho_\tau \sim 10^{20}\text{--}10^{25}\text{ Hz}\)) to resolve the
flux, while leaving structures too rigid for historical memory
(\(K \to 0\)). * \textbf{The Mesoscopic Biological Complexity Peak
(\(10^{-9}\text{--}10^{0}\text{ m}, T_{\text{coh}} \approx 300\text{ K}\)):}
At biological temperatures,
\(\mathcal{L}_{\text{sim}}(300\text{ K}) = \hbar c / (2\pi k_B T_{\text{coh}}) \approx 1.21\,\mu\text{m}\)---the
exact scale of a living cell. Sitting squarely on the knife-edge where
\(\Delta x \approx \mathcal{L}_{\text{sim}}\), living systems harvest
quantum phase-coherence (\(\Omega\)) while using selective thermal
condensation (\(\Phi_\tau\)) to write irreversible, low-entropy
historical records (\(K_{\text{max}}\): DNA, proteins, neural memory).

\begin{center}\rule{0.5\linewidth}{0.5pt}\end{center}

\subsection{Appendix: The Dual Law of Chronometric Reality --- Unifying
the Cadence Partition Law and the Horizon-Condensation
Law}\label{appendix-the-dual-law-of-chronometric-reality-unifying-the-cadence-partition-law-and-the-horizon-condensation-law}

\subsubsection{A.1 Intra-Entity Allocation and Inter-Entity
Synchronization}\label{a.1-intra-entity-allocation-and-inter-entity-synchronization}

In the Theory of Objective Time (\(\text{AEON-}T_{\mu\nu}\)), physical
existence is governed by two complementary, mutually coupling laws that
regulate \textbf{Temporal Intensity (\(\rho_\tau\))} across individual
entities and across extended spatial horizons:

\begin{enumerate}
\def\labelenumi{\arabic{enumi}.}
\item
  \textbf{The Cadence Partition Law (Intra-Entity / Vertical Axis):}
  \[\rho_{\tau,\text{total}}^2 = \rho_{\tau,\text{int}}^2 + \rho_{\tau,\text{kin}}^2\]
  Governs how a single existent allocates its finite budget of
  ontological resolution between \textbf{immanent self-maintenance}
  (\(\rho_{\tau,\text{int}}\), chiral phase-locked rest mass) and
  \textbf{transitive spatial displacement} (\(\rho_{\tau,\text{kin}}\),
  relational momentum), parameterized by the partition angle \(\theta\)
  (\(\sin\theta = \beta, \cos\theta = 1/\gamma\)).
\item
  \textbf{The Horizon-Condensation Law (Inter-Entity / Horizontal
  Axis):}
  \[\mathcal{C}_\tau = \frac{k_B T_{\text{coh}} \cdot \mathcal{L}_{\text{sim}}}{h c} = \frac{1}{2\pi} \quad \Longleftrightarrow \quad \frac{k_B T_{\text{coh}} \cdot \mathcal{L}_{\text{sim}}}{\hbar c} = 1\text{ rad}\]
  Governs the maximum spatial radius \(\mathcal{L}_{\text{sim}}\) over
  which distinct systems can maintain coherent phase-lock inside a
  unified \textbf{Functional Infinity (\(\Omega\))} before stochastic
  dispersion triggers \textbf{Two-Stage Condensation
  (\(\Phi_\tau, v \to v + 1\))}.
\end{enumerate}

\subsubsection{\texorpdfstring{A.2 Geometric Derivation of
\(\mathcal{C}_\tau = 1/2\pi\) from Partition-Angle Jitter
(\(\delta\theta\))}{A.2 Geometric Derivation of \textbackslash mathcal\{C\}\_\textbackslash tau = 1/2\textbackslash pi from Partition-Angle Jitter (\textbackslash delta\textbackslash theta)}}\label{a.2-geometric-derivation-of-mathcalc_tau-12pi-from-partition-angle-jitter-deltatheta}

The dimensionless constant \(\mathcal{C}_\tau = 1/2\pi\) is not an
arbitrary numerical parameter; it is the \textbf{1-radian phase-slip
threshold} linking partition-angle fluctuations (\(\delta\theta\)) to
spatial horizon collapse:

\begin{itemize}
\tightlist
\item
  When an ensemble exhibits an effective phase-coherence temperature
  \(T_{\text{coh}}\), random fluctuations \(\delta\theta\) in each
  entity's Cadence Partition angle generate a stochastic dispersion in
  transitive angular frequency:
  \[\delta\omega_{\text{kin}} = \omega_{\tau,\text{int}}\,\delta\theta = \frac{k_B T_{\text{coh}}}{\hbar}.\]
\item
  Across a spatial separation \(\Delta x\) bridged by light-speed
  substrate coupling (\(\Delta t_{\text{sig}} = \Delta x / c\)), this
  dispersion accumulates a relative phase drift:
  \[\Delta\phi_{\text{slip}}(\Delta x) = \delta\omega_{\text{kin}} \cdot \frac{\Delta x}{c} = \frac{k_B T_{\text{coh}} \cdot \Delta x}{\hbar c}.\]
\item
  Coherent phase-locking inside \(\Omega\) persists strictly while
  \(\Delta\phi_{\text{slip}} < 1\text{ rad}\) (less than \(1/2\pi\) of a
  complete \(2\pi\)-radian standing cycle). Setting
  \(\Delta\phi_{\text{slip}}(\mathcal{L}_{\text{sim}}) = 1\text{ rad}\)
  yields the Weyl-invariant horizon constant:
  \[\mathcal{C}_\tau = \frac{k_B T_{\text{coh}} \cdot \mathcal{L}_{\text{sim}}}{h c} = \frac{1}{2\pi}.\]
\end{itemize}

\subsubsection{A.3 The Closed Causal Loop Between Both
Laws}\label{a.3-the-closed-causal-loop-between-both-laws}

Neither law operates in isolation; together they form a closed
ontological feedback loop:

{\def\LTcaptype{none} % do not increment counter
\begin{longtable}[]{@{}
  >{\raggedright\arraybackslash}p{(\linewidth - 4\tabcolsep) * \real{0.3333}}
  >{\raggedright\arraybackslash}p{(\linewidth - 4\tabcolsep) * \real{0.3333}}
  >{\raggedright\arraybackslash}p{(\linewidth - 4\tabcolsep) * \real{0.3333}}@{}}
\toprule\noalign{}
\begin{minipage}[b]{\linewidth}\raggedright
Transition
\end{minipage} & \begin{minipage}[b]{\linewidth}\raggedright
Mechanism
\end{minipage} & \begin{minipage}[b]{\linewidth}\raggedright
Physical \& Ontological Consequence
\end{minipage} \\
\midrule\noalign{}
\endhead
\bottomrule\noalign{}
\endlastfoot
\textbf{Partition \(\to\) Horizon}
(\(\rho_{\tau,\text{int}} \to \mathcal{L}_{\text{sim}}\)) & Pure
transitive modes (\(\rho_{\tau,\text{int}} = 0, T_{\text{coh}} \to 0\))
have zero partition jitter (\(\delta\theta = 0\)), giving
\(\mathcal{L}_{\text{sim}} \to \infty\). Once chiral locking
(\(\psi_L \leftrightarrow \psi_R\)) creates massive clocks
(\(\rho_{\tau,\text{int}} > 0\)), collisional exchanges inject jitter
\(T_{\text{coh}} > 0\), bounding \(\mathcal{L}_{\text{sim}}\) to a
finite radius. & Explains why massless photons span cosmic horizons
un-condensed while thermalized matter condenses into localized classical
objects (\(1.21\,\mu\text{m}\) at \(300\text{ K}\)). \\
\textbf{Horizon \(\to\) Partition} (\(\Phi_\tau \to \theta + d\theta\))
& Whenever two systems separate beyond
\(\Delta x \ge \mathcal{L}_{\text{sim}}(T_{\text{coh}})\), phase-slip
exceeds \(1\text{ rad}\) and triggers \(\Phi_\tau\) condensation
(\(v \to v + 1\)), exchanging discrete quanta of
\(\rho_{\tau,\text{kin}}\) and rotating each entity's partition angle
\(\theta\). & Explains how discrete interactions, mechanical forces
(\(F = dp/dt\)), and \emph{in situ} clock measurements
(\emph{object-as-subject}) emerge at the boundary of
\(\mathcal{L}_{\text{sim}}\). \\
\end{longtable}
}

\begin{center}\rule{0.5\linewidth}{0.5pt}\end{center}

\subsection{Bibliography \& Reference
Framework}\label{bibliography-reference-framework}

\begin{enumerate}
\def\labelenumi{\arabic{enumi}.}
\tightlist
\item
  \textbf{Pimenta Damásio, L.} (2026). \emph{Master Compendium:
  Foundations of Temporal Metrology and Operational Quantum Mechanics}
  (\(\text{AEON-}T_{\mu\nu}\)). Zenodo. DOI:
  \texttt{10.5281/zenodo.24163322}. ORCID: \texttt{0009-0001-7606-7557}.
\item
  \textbf{Pimenta Damásio, L.} (2026). \emph{Coherence-Bounded
  Operational Locality: An Experimental Test of Simultaneous Phase
  Synchronization vs.~Bell Correlations} (Paper 10).
\item
  \textbf{Kant, I.} (1781 / 1787). \emph{Critique of Pure Reason}
  (\emph{Transcendental Aesthetic}).
\item
  \textbf{Einstein, A.} (1905). ``On the Electrodynamics of Moving
  Bodies'' \& ``Does the Inertia of a Body Depend Upon Its
  Energy-Content?'' \emph{Annalen der Physik}, 17 \& 18.
\item
  \textbf{Planck, M.} (1901). ``On the Law of Distribution of Energy in
  the Normal Spectrum.'' \emph{Annalen der Physik}, 4, 553--563.
\item
  \textbf{Nambu, Y., \& Jona-Lasinio, G.} (1961). ``Dynamical Model of
  Elementary Particles Based on an Analogy with Superconductivity.''
  \emph{Physical Review}, 122(1), 345--358.
\item
  \textbf{Mannheim, P. D.} (2012). ``Making the Case for Conformal
  Gravity.'' \emph{Foundations of Physics}, 42(3), 388--420.
\item
  \textbf{Pimenta Dam\'asio, L.} (2026). ``The Horizon-Condensation Law
  ($\mathcal{C}_\tau = 1/2\pi$): Speed, Temperature, and the Weyl-Invariant
  Bounds of Operational Simultaneity.'' \emph{Zenodo}.
  \url{https://doi.org/10.5281/zenodo.24444831}.
\item
  \textbf{Pimenta Dam\'asio, L.} (2026). ``AEON-T$\mu\nu$: Covariant
  Kerr--Newman QFTCS \& Semiclassical Back-Reaction Observatory'' [Computer
  software]. \emph{Zenodo}. \url{https://doi.org/10.5281/zenodo.24244501}.
\end{enumerate}

\begin{center}\rule{0.5\linewidth}{0.5pt}\end{center}

\subsection{Author \& AI Assistance
Statement}\label{author-ai-assistance-statement}

\begin{itemize}
\tightlist
\item
  \textbf{Conceptual \& Ontological Authorship:} All foundational
  concepts, philosophical postulates, the Theory of Objective Time
  (\(\text{AEON-}T_{\mu\nu}\)), the Functional Infinity (\(\Omega\)),
  Two-Stage Condensation (\(\Phi_\tau\)), the Cadence Partition Law, and
  the Horizon-Condensation Law are the original work of \textbf{Luís
  Pimenta Damásio} (ORCID: \texttt{0009-0001-7606-7557}).
\item
  \textbf{AI Assistance Disclaimer:} This manuscript was structured,
  formatted, and mathematically cross-checked with the assistance of an
  AI language model acting as an editorial and formalization aid under
  the author's direct conceptual direction.
\end{itemize}

\end{document}
